\begin{center}
\LARGE
Symbolic Methods in Analog Circuit Design
\end{center}

\begin{center}
\Large
Ralf Sommer, Eckhard Hennig

\end{center}

\noindent
Date:  July 19th (Friday)\\
Time:  08:30-09:00\\

\begin{center}
\large
Abstract
\end{center}

This presentation will give an introduction to the computer-aided approach to the analysis and
design of analog electronic circuits. At the beginning, typical engineering tasks will be
characterized by several examples, and their mathematical background will be formulated. This
will include different techniques for setting up circuit equations. Hence, network elements and
their characteristics as well as Kirchhoff laws will be covered. Then the symbolic formulation of
analysis and design equations will be looked into. 

Special attention will be attached to the problem of solving the equations since the explosion of
terms in the solutions is the main problem in symbolic circuit analysis. Three fundamental
approaches for the automated approximation of circuit transfer function and circuit equations will
be presented which allow for the extraction of dominant circuit behaviour - linear and even
nonlinear. These approximation strategies are based on combining symbolic and numeric
information: They hold the key for gaining insight in how circuits work, and the obtained results
may even be used for the generation of sizing formulas. From a mathematical point of view these
new mixed symbolic-numerical approximation strategies open up a large field of unsolved
problems. 

